% Part: clearly prospective unification; no BFSS particle prediction claimed.
\chapter[Representation Structure and Particle Families]{Prospective Particle Structure: Spin Representations, Gauge Color, and the Three-Family Test}
\label{ch:gut-representations}

\section{What may be asked without changing the quantum model}
The source contains a real nine-dimensional spatial vector
representation, a real sixteen-dimensional Spin(9) spinor, and
three real color components in the adjoint representation of SU(2).
Thus at the level of Majorana \emph{labels}, after the exact
reindexing established in Chapter~\ref{ch:spin-geometry},
\begin{equation}\label{eq:unification-labels}
 V_{48}\cong S_{16,\mathbb R}\otimes\mathfrak{su}(2)_{\mathbb R},
 \qquad \dim_{\mathbb R}V_{48}=16\cdot3=48.
\end{equation}
This factorization is a powerful source-faithful organizing device.
It is \emph{not} three gauge-invariant fermion generations. The
SU(2) adjoint group acts nontrivially on the entire color factor,
while the source Spin(9) group acts on the spinor factor. The
observed Standard Model generations, if they are to emerge, require
physically separate chiral particle sectors, not merely three
basis directions in an internal gauge representation.

\section{A precisely limited color-projector obstruction}
\begin{proposition}[No three invariant rank-one color labels;
conventional representation-theory result]\label{prop:no-three-color-projectors}
Let $R:SU(2)\to SO(3)$ be its real adjoint rotation
representation. If a real linear map $P:\mathbb R^3\to\mathbb R^3$
commutes with $R(g)$ for every $g\in SU(2)$, then $P=cI_3$ for
some $c\in\mathbb R$. In particular there cannot be three distinct
nonzero rank-one projections individually commuting with every
source color gauge transformation.
\end{proposition}
\begin{proof}
The image of the adjoint representation is the full $SO(3)$
rotation group. For any unit vector $v$, rotations fixing $v$
force $Pv$ to lie in the line through $v$: the only vectors fixed
by all rotations about $v$ are its multiples. Hence every nonzero
$v$ is an eigenvector of $P$. Applying this to $u$, $v$, and $u+v$
for linearly independent $u,v$ shows that the eigenvalues are equal.
Thus $P$ is scalar. A nonzero scalar idempotent has rank three,
not one.
\end{proof}
\formalrecord{The source-specific coefficients and the original
G3A/C48-0B identification are accepted; the explicit common
commutant calculation in
\path{research/unification/BFSS_THREE_BY_SIXTEEN_SPIN9_SPIN10_FAMILY_FEASIBILITY_2026_11_10.md}
was a conventional/symbolic audit, \textbf{not a new original-source Lean
commutant theorem}. This proposition rules out only the naive
``one gauge color equals one generation'' mechanism.}

\section{Why the Spin(10) sixteen is a lead, not a conclusion}
The standard inclusion $\operatorname{Spin}(9)\subset
\operatorname{Spin}(10)$ has the familiar representation-theoretic
property that either chiral complex sixteen-dimensional spinor of
Spin(10) restricts to the complexification of the real Spin(9)
sixteen. This gives a legitimate abstract mathematical comparison.
But the two chiral Spin(10) representations become indistinguishable
under this restriction, and the source BFSS spinor index is a
\emph{spatial rotational} spinor, whereas a putative Spin(10) GUT
multiplet would be an \emph{internal gauge} representation of
four-dimensional chiral matter. Converting one into the other
requires an explicit physical mechanism and new, source-bound
structures. Numerical matching of dimensions is insufficient.

Even the Clifford-algebra statement
$\operatorname{Cl}(V\oplus W)\cong
\operatorname{Cl}(V)\widehat\otimes\operatorname{Cl}(W)$
is a \emph{graded} factorization; it does not decompose the full
Spin(48) group into three independent Spin(16) factors or produce
three physical superselection sectors. The cross-block bivectors
remain part of the full $\mathfrak{so}(48)$ algebra. These are
important safeguards against a speculative but invalid family count.

\section{A falsifiable candidate-family protocol}
Any new three-generation candidate must produce actual gauge-invariant
physical subspaces, a well-defined chiral four-dimensional
realization, correct Standard Model hypercharges and nonabelian
representations, anomaly-consistent matter content, and a dynamics
that makes the proposed families observable. It should also explain
why the multiplicity is exactly three rather than one, two, or any
chosen number. A mathematical construction that fails any essential
condition is to be preserved as a no-go or a narrower algebraic
result, not called a partially verified GUT.

We therefore give priority to (i) a source-typed color commutant/no-go
theorem, (ii) the original Spin(9) $16\to8+8$ eigenspace
restriction under its actual axis stabilizer, (iii) a careful
abstract Spin(10) branching crosswalk, and (iv) gauge-invariant
higher-degree Fock composites. None of these is currently a
three-generation physical theorem.
